\lecture{5}{2026-09-18}{Analytic continuation and complex integration}

\section*{Analytic continuation}

By the Identity Theorem, an analytic function on a domain \(\Omega\) is uniquely determined by its values on any non-empty open subset of \(\Omega\). This property provides a principled way to extend the domain of definition of analytic functions.

If \(\Omega_1\) and \(\Omega_2\) are two domains in \(\mathbb{C}\) with non-empty intersection \(\Omega_1 \cap \Omega_2 \neq \varnothing\), and 
\(f \colon \Omega_1 \to \mathbb{C}\) and \(g \colon \Omega_2 \to \mathbb{C}\) are analytic functions that agree on \(\Omega_1 \cap \Omega_2\), then the function \(h \colon \Omega_1 \cup \Omega_2 \to \mathbb{C}\) defined by
\[
    h(z) = \begin{cases}
        f(z) & z \in \Omega_1, \\
        g(z) & z \in \Omega_2,
    \end{cases}
\]
is well-defined and analytic on the domain \(\Omega_1 \cup \Omega_2\).

In this case, \(g\) is called a \emph{direct analytic continuation} of \(f\) to \(\Omega_2\), and \(h\) is the analytic continuation of \(f\) and \(g\) to \(\Omega_1 \cup \Omega_2\).
More generally, we say that \(g\) is an analytic continuation of \(f\) if there exists a chain of direct analytic continuations on consecutively overlapping domains \(\Omega_1, \Omega_2, \dots, \Omega_r\) beginning with \(f\) on \(\Omega_1\) and ending with \(g\) on \(\Omega_r\).

\begin{example}[Analytic continuation of a geometric series]
    The rational function \(\frac{1}{1-z}\) is a direct analytic continuation of the power series \(\sum_{n \ge 0} z^n\) from the unit disk \(D_0^{(1)}\) to the domain \(\mathbb{C} \setminus \{1\}\).
\end{example}

Note that general analytic continuations are not unique.

\pagebreak

\begin{example}[Non-uniqueness of analytic continuation]
    Consider the right half-plane \(\Omega_0 = \{ z \in \mathbb{C} : \Re(z) > 0 \}\), where the square root function is defined by
    \[
        F_0(z) = \sqrt{r} e^{i\theta/2} \quad \text{for } z = r e^{i\theta} \text{ with } r > 0 \text{ and } \theta \in (-\pi/2, \pi/2),
    \]
    so that \(F_0(1) = 1\).
    We can analytically continue \(F_0\) to domains containing the point \(z = -1\) in different ways.

    Let \(\Omega_1 = \mathbb{C} \setminus \{ -iy : y \ge 0 \}\) (the complex plane excluding the non-positive imaginary ray).
    Every \(z \in \Omega_1\) has polar angle \(\theta \in (-\pi/2, 3\pi/2)\).
    Defining \(F_1(r e^{i\theta}) = \sqrt{r} e^{i\theta/2}\), we have \(F_1(z) = F_0(z)\) on the overlap \(\Omega_0 \cap \Omega_1 = \Omega_0\).
    Hence, \(F_1\) is a direct analytic continuation of \(F_0\) to \(\Omega_1\), and since \(-1 = 1 \cdot e^{i\pi} \in \Omega_1\),
    \( F_1(-1) = e^{i\pi/2} = i \).

    \begin{figure}[b]
        \begin{minipage}[t]{0.48\textwidth}
            \centering
            \begin{tikzpicture}[scale=0.85, >=Stealth]
                % Domain Omega_1: fill entire displayed region except negative imaginary axis
                \fill[blue!6, rounded corners=3pt] (-2.2,-2.0) rectangle (2.2,2.0);
                
                % Highlight initial domain Omega_0 (right half plane)
                \fill[blue!14] (0,-2.0) rectangle (2.2,2.0);

                \draw[ line cap=butt, draw=white, line width=3pt, shorten < = -1.5pt, {Round Cap[sep=-0.5]}- ] (0,0) -- (0, -2.0);
                
                % Axes
                \draw[thick, ->] (-2.3, 0) -- (2.3, 0) node[right] {\footnotesize $\Re(z)$};
                \draw[thick, ->] (0, -2.1) -- (0, 2.4) node[right] {\footnotesize $\Im(z)$};

                \draw[draw=blue!80!black, line width=1pt, line cap=butt] (2.0pt, -2.0) -- (2.0pt, 0) arc[start angle=0, end angle=180, radius=2.0pt] -- (-2.0pt, -2.0);
                
                % Start point z = 1
                \filldraw[black] (1.4,0) circle (2pt);
                \node[below=2pt, font=\footnotesize] at (1.4,0) {$1$};
                \node[below=2pt, blue!85!black, font=\tiny] at (1.2, -0.6) {$F_0(1)=1$};
                
                % Target point z = -1
                \filldraw[black] (-1.4,0) circle (2pt);
                \node[below=2pt, font=\footnotesize] at (-1.4,0) {$-1$};
                \node[below=2pt, blue!85!black, font=\tiny] at (-1.2, -0.6) {$F_1(-1) = +i$};
                
                % Arc via top (argument goes from 0 to pi)
                \draw[very thick, blue!80!black, ->] (1.4,0) arc (0:95:1.4);
                \draw[very thick, blue!80!black] (0,1.4) arc (90:180:1.4);
            \end{tikzpicture}
            \par\vspace{0.4em}
            {\footnotesize \textbf{(a)} Continuation \(F_1\) to \(\Omega_1\)}
        \end{minipage}\hfill
        \begin{minipage}[t]{0.48\textwidth}
            \centering
            \begin{tikzpicture}[scale=0.85, >=Stealth]
                % Domain Omega_2: fill entire displayed region except positive imaginary axis
                \fill[teal!6, rounded corners=3pt] (-2.2,-2.0) rectangle (2.2,2.0);
                
                % Highlight initial domain Omega_0 (right half plane)
                \fill[teal!14] (0,-2.0) rectangle (2.2,2.0);

                \draw[ line cap=butt, draw=white, line width=3pt, shorten < = -1.5pt, {Round Cap[sep=-0.5]}- ] (0,0) -- (0, 2.0);
                
                % Axes
                \draw[thick, ->] (-2.3, 0) -- (2.3, 0) node[right] {\footnotesize $\Re(z)$};
                \draw[thick, ->] (0, -2.1) -- (0, 2.4) node[right] {\footnotesize $\Im(z)$};

                \draw[draw=teal!80!black, line width=1pt, line cap=butt] (2.0pt, 2.0) -- (2.0pt, 0) arc[start angle=0, end angle=-180, radius=2.0pt] -- (-2.0pt, 2.0);
                
                % Start point z = 1
                \filldraw[black] (1.4,0) circle (2pt);
                \node[above=2pt, font=\footnotesize] at (1.4,0) {$1$};
                \node[above=2pt, teal!85!black, font=\tiny] at (1.2, 0.6) {$F_0(1)=1$};
                
                % Target point z = -1
                \filldraw[black] (-1.4,0) circle (2pt);
                \node[above=2pt, font=\footnotesize] at (-1.4,0) {$-1$};
                \node[above=2pt, teal!85!black, font=\tiny] at (-1.2, 0.6) {$F_2(-1) = -i$};
                
                % Arc via bottom (argument goes from 0 to -pi)
                \draw[very thick, teal!80!black, ->] (1.4,0) arc (0:-95:1.4);
                \draw[very thick, teal!80!black] (0,-1.4) arc (-90:-180:1.4);
            \end{tikzpicture}
            \par\vspace{0.4em}
            {\footnotesize \textbf{(b)} Continuation \(F_2\) to \(\Omega_2\)}
        \end{minipage}
        \caption{Two distinct direct analytic continuations of \(F_0(z) = \sqrt{r}e^{i\theta/2}\).}
    \end{figure}

    Let \(\Omega_2 = \mathbb{C} \setminus \{ iy : y \ge 0 \}\), with polar angles \(\theta \in (-3\pi/2, \pi/2)\), and define
    \( F_2(r e^{i\theta}) = \sqrt{r}e^{i\theta/2} \).
    This agrees with \(F_0\) on \(\Omega_0\), so it is a direct continuation. Since \(-1=e^{-i\pi}\in\Omega_2\),
    \( F_2(-1) = e^{-i\pi/2} = -i \).

    Thus both continuations contain \(-1\), but
    \[
        F_1(-1) = i \neq -i = F_2(-1).
    \]
    Therefore, continuation to a distant point can depend on the domains chosen, even though continuation to a fixed overlapping domain is unique by the Identity Theorem.
\end{example}


\section*{Complex integration}

In real analysis, we integrate over intervals \([a, b] \subset \mathbb{R}\).
In complex analysis, we integrate over curves in \(\mathbb{C}\).

A \emph{curve} is a piecewise continuously differentiable function \(\gamma \colon [a, b] \to \mathbb{C}\).
Writing \(\gamma(t) = x(t) + i y(t)\), the components \(x(t), y(t) \colon [a, b] \to \mathbb{R}\) are piecewise continuously differentiable real functions.

\begin{definition}[Contour integral]
If \(f\) is continuous on the image \(\gamma([a, b])\), the \emph{contour integral} of \(f\) along \(\gamma\) is defined by
\[
    \int_\gamma f(z) \, dz = \int_a^b f(\gamma(t)) \gamma'(t) \, dt.
\]
\end{definition}

If \(f(z) = u(z) + i v(z)\) with \(u, v \colon \mathbb{C} \to \mathbb{R}\), linearity implies
\[
    \int_\gamma f(z) \, dz = \int_\gamma u(z) \, dz + i \int_\gamma v(z) \, dz.
\]

\begin{proposition}
    The contour integral depends only on the oriented trace of the curve, not on the parametrization.
    Formally, if \(\gamma_1 \colon [a, b] \to \mathbb{C}\) and \(\gamma_2 \colon [c, d] \to \mathbb{C}\) are two curves and \(h \colon [a, b] \to [c, d]\) is an orientation-preserving continuously differentiable bijection (with \(h'(t) > 0\) and continuously differentiable inverse) satisfying \(\gamma_1(t) = \gamma_2(h(t))\) for all \(t \in [a, b]\), then
    \[
        \int_{\gamma_1} f(z) \, dz = \int_{\gamma_2} f(z) \, dz.
    \]
\end{proposition}

\begin{example}[Contour integrals of \(z^n\)]
    Let \(C\) be the unit circle \(\{|z| = 1\}\) oriented counter-clockwise.
    We compute the integral
    \[
        \int_C z^n \, dz
    \]
    for each \(n \in \mathbb{Z}\).
    Parametrizing \(C\) by \(\gamma(t) = e^{it}\) for \(t \in [-\pi, \pi]\), so that \(\gamma'(t) = i e^{it}\):
    
    If \(n \in \mathbb{Z} \setminus \{-1\}\), then
    \begin{align*}
        \int_C z^n \, dz
        &= \int_{-\pi}^{\pi} \gamma(t)^n \gamma'(t) \, dt \\
        &= \int_{-\pi}^{\pi} e^{int} i e^{it} \, dt \\
        &= i \int_{-\pi}^{\pi} e^{i(n+1)t} \, dt \\
        &= \left[ \frac{e^{i(n+1)t}}{n+1} \right]_{-\pi}^{\pi} \\
        &= 0,
    \end{align*}
    since \(e^{i(n+1)\pi} = e^{-i(n+1)\pi} = (-1)^{n+1}\).

    If \(n = -1\), then
    \begin{align*}
        \int_C z^{-1} \, dz
        &= \int_{-\pi}^{\pi} \gamma(t)^{-1} \gamma'(t) \, dt \\
        &= \int_{-\pi}^{\pi} e^{-it} i e^{it} \, dt \\
        &= i \int_{-\pi}^{\pi} 1 \, dt \\
        &= 2\pi i.
    \end{align*}
    This calculation is crucial for analytic combinatorics: the integral picks out precisely the \(z^{-1}\) term and eliminates all others, making contour integration the primary tool for coefficient extraction.
\end{example}

\subsection*{Topology of curves and Cauchy's theorem}

A curve \(\gamma \colon [a, b] \to \mathbb{C}\) is:
\begin{itemize}
    \item \emph{closed} if \(\gamma(a) = \gamma(b)\);
    \item \emph{simple} if \(\gamma\) is injective on \([a, b)\); that is, it does not self-intersect (except possibly at the endpoints \(\gamma(a) = \gamma(b)\)).
\end{itemize}

Two curves \(\gamma_1, \gamma_2\) in a domain \(\Omega\) sharing the same endpoints are called \emph{homotopic} in \(\Omega\) if one can be continuously deformed into the other while remaining entirely inside \(\Omega\).

Following the terminology in \cite{Melczer2021}, a \emph{loop} in a domain \(\Omega\) is a simple closed curve that can be continuously deformed to a single point while staying in \(\Omega\) (i.e., a contractible or null-homotopic simple closed curve).
A domain \(\Omega\) is \emph{simply connected} if every simple closed curve in \(\Omega\) is a loop in \(\Omega\).

For instance, in the punctured plane \(\Omega = \mathbb{C} \setminus \{0\}\), the circle \(\gamma(t) = e^{it}\) (\(t \in [0, 2\pi]\)) is a simple closed curve, but it cannot be contracted to a point within \(\Omega\) because the puncture at the origin obstructs the deformation. Thus, \(\mathbb{C} \setminus \{0\}\) is not simply connected.

\begin{theorem}[Jordan curve theorem]
    Let \(\gamma\) be a simple closed curve in \(\mathbb{C}\).
    Then \(\mathbb{C} \setminus \gamma\) consists of exactly two connected components: one bounded (called \emph{the interior}) and one unbounded (called \emph{the exterior}), with \(\gamma\) being the common boundary of both components.
\end{theorem}

A simple closed curve is \emph{positively oriented} if its interior lies to the left as the curve is traversed.
For example, the circle \(\gamma(t) = e^{it}\) for \(t \in [-\pi, \pi]\) is positively oriented (counter-clockwise).

\begin{proposition}[Path independence and Cauchy's theorem]
    If \(f\) is analytic in a domain \(\Omega\) and \(\gamma_1, \gamma_2\) are two curves in \(\Omega\) with the same endpoints that are homotopic in \(\Omega\), then
    \[
        \int_{\gamma_1} f(z) \, dz = \int_{\gamma_2} f(z) \, dz.
    \]
    In particular, if \(\gamma\) is a loop in \(\Omega\), then
    \[
        \int_\gamma f(z) \, dz = 0.
    \]
\end{proposition}

This implies that contour integrals of analytic functions depend only on the homotopy class of the path; contours can be continuously deformed across regions of analyticity without altering the value of the integral.

\begin{proposition}[Maximum modulus integral bound (\(ML\)-inequality)]
    If \(f\) is continuous on a curve \(\gamma\), then
    \[
        \left| \int_\gamma f(z) \, dz \right| \le \operatorname{length}(\gamma) \cdot \max_{z \in \gamma} |f(z)|.
    \]
\end{proposition}

\begin{proof}
    Applying the triangle inequality for integrals:
    \begin{align*}
        \left| \int_\gamma f(z) \, dz \right|
        &= \left| \int_a^b f(\gamma(t)) \gamma'(t) \, dt \right| \\
        &\le \int_a^b |f(\gamma(t))| |\gamma'(t)| \, dt \\
        &\le \left( \max_{z \in \gamma} |f(z)| \right) \int_a^b |\gamma'(t)| \, dt = \operatorname{length}(\gamma) \cdot \max_{z \in \gamma} |f(z)|. \qedhere
    \end{align*}
\end{proof}